\slcol{अर्जुन उवाच 
\Index{एवं सततयुक्ता} ये भक्तास्त्वां पर्युपासते ।\\
ये चाप्यक्षरमव्यक्तं तेषां के योगवित्तमाः ॥ १२.१ ॥}
\translate{arjuna uvāca—\\
evaṁ satata-yuktā ye bhaktās tvāṁ paryupāsate |\\
ye cāpy akṣaram avyaktam teṣāṁ ke yoga-vittamāḥ || 12.1 ||}
\cquote{Arjuna said: Of those devotees who, ever steadfast, worship You with devotion, and those who worship the imperishable, the unmanifest (\emph{avyaktam})—which of them are better versed in \emph{yoga}?}
\slcol{श्रीभगवानुवाच \\
\Index{मय्यावेश्य मनो} ये मां नित्ययुक्ता उपासते ।\\
श्रद्धया परयोपेताः ते मे युक्ततमा मताः ॥ १२.२ ॥}
\translate{śrī-bhagavān uvāca\\
mayy āveśya mano ye māṁ nitya-yuktā upāsate |\\
śraddhayā parayopetāḥ te me yukta-tamā matāḥ || 12.2 ||}
\cquote{Sri Bhagavān said: Those who, fixing their mind on Me, worship Me, ever steadfast and endowed with supreme faith, these, in my opinion, are the best in \emph{yoga}.}

\newpage
\begin{mananam}{\mananamfont {Mananam Śloka - 1}}
\normalfont\mananamtext Which concept of the Divine resonates with me the most? Do I find joy in contemplating a personal form of God? Do stories of the personal God such as those in the Indian \emph{Puranas} uplift my mind and inspire hope?
Or, do I perhaps intuitively connect with a formless divine force—the unmanifest energy behind all material existence? Am I able to embrace both perspectives in my understanding and relate to them based on my mental state and emotional needs?
\end{mananam}
\sreflect
\begin{inspiration}{\inspirationTitleFont{Inspiration}}
\normalfont\mananamtext The divine force behind creation and consciousness has been understood in different ways by seekers. Among them, the majority embrace the worship of God in a specific form with distinct attributes. In contrast, \emph{yogīs} and \emph{jñānīs} transcend the need for a defined form, seeking instead the essence of the Divine within their own hearts—the one truth behind all appearances. 
In essence, both the personal and impersonal aspects of God represent the same divine reality, approached in different ways according to an individual's temperament and state of mind.
\end{inspiration}
\newpage

\newpage
\begin{mananam}{\mananamfont {Mananam Śloka - 2}}
\normalfont\mananamtext In what ways can devotion to a particular aspect of God be fulfilling? How does faith in the Divine help me to navigate the challenges of everyday life? Is faith merely a source of emotional support, or does it bring tangible benefits to my life experiences? How can I train my mind to remain steadfastly focused on the Lord? What steps can I take to cultivate deeper devotion and unwavering faith?
\end{mananam}
\sreflect
\begin{inspiration}{\inspirationTitleFont{Inspiration}}
\normalfont\mananamtext
The best kind of faith is one which arises from personal understanding and experience. However, it is every devotee’s responsibility to actively seek the Lord, following the guidance of the scriptures and \emph{gurus}. Through sincere effort and devotion, this faith gradually matures until it manifests fully in one's life.\\
A true devotee finds ultimate refuge in the Lord, recognizing that all other sources—family, power, wealth—are secondary and merely manifestations of divine grace. As this perspective deepens, life itself becomes a sacred journey of anchoring oneself ever more firmly in the Lord. 
\end{inspiration}
\newpage

\slcol{\Index{ये त्वक्षरमनिर्देश्यमव्यक्तं} पर्युपासते ।\\
सर्वत्रगमचिन्त्यं च कूटस्थमचलं ध्रुवम् ॥ १२.३ ॥}
\translate{ye tv akṣaram anirdeśyam avyaktam paryupāsate |\\
sarvatra-gam acintyaṁ ca kūṭastham acalaṁ dhruvam || 12.3 ||}
\cquote{But those who worship the imperishable, the indefinable, the unmanifest, the all-pervading, the inconceivable, the unchanging, the immovable, and the eternal…}
\slcol{\Index{संयम्येन्द्रियग्रामं} सर्वत्र समबुद्धयः ।\\
ते प्राप्नुवन्ति मामेव सर्वभूतहिते रताः ॥ १२.४ ॥}
\translate{saṁyamya indriya-grāmaṁ sarvatra sama-buddhayaḥ |\\
te prāpnuvanti mām eva sarva-bhūta-hite ratāḥ || 12.4 ||}
\cquote{Controlling all the senses, being even-minded everywhere, and engaged in the welfare of all beings—they attain Me alone.}
\slcol{\Index{क्लेशोऽधिकतरस्तेषा}मव्यक्तासक्तचेतसाम् ।\\
अव्यक्ता हि गतिरदुःखं देहवद्भिरवाप्यते ॥ १२.५ ॥}
\translate{kleśo ’dhikataras teṣām avyaktāsakta-cetasām |\\
avyaktā hi gatir duḥkhaṁ dehavadbhir avāpyate || 12.5 ||}
\cquote{For those whose minds are attached to the \emph{avyaktā}  (unmanifest), greater is the struggle. Attaining the goal of \emph{avyaktā} is fraught with hardship for the embodied ones.}
\slcol{\Index{ये तु सर्वाणि} कर्माणि मयि संन्यस्य मत्पराः ।\\
अनन्येनैव योगेन मां ध्यायन्त उपासते ॥ १२.६ ॥}
\translate{ye tu sarvāṇi karmāṇi mayi saṁnyasya mat-parāḥ |\\
ananyenaiva yogena māṁ dhyāyanta upāsate || 12.6 || }
\cquote{But those who, dedicating all their actions to Me, considering Me as their supreme goal, worship me, meditate upon Me with undivided devotion…. }

\newpage
\begin{mananam}{\mananamfont {Mananam Śloka - 3, 4}}
\normalfont\mananamtext How may I grasp the formless aspect of Divinity, which is beyond any thought? Or rather, do I dismiss anything beyond human comprehension as non-existent?
Have I ever tried to experience God as a silent presence, a sudden spark of peace and joy within my heart, or a surge of a deep sense of goodwill toward all humanity?
What inner conditions can help me cultivate awareness of the unmanifest yet omnipresent aspect of the Divine?
\end{mananam}
\sreflect
\begin{inspiration}{\inspirationTitleFont{Inspiration}}
\normalfont\mananamtext Devotion to a personal God is often shaped by cultural and social influences, but the experience of the formless, impersonal Divine can only be attained as a result of dedicated \emph{sādhana}, such as meditation and other spiritual practices. When a \emph{mumukṣu} (seeker) with a purified and prepared mind receives the wisdom of the scriptures through a \emph{guru}, it leads toward the direct realisation of the eternal and transcendental nature of God. Upon awakening to this unmanifest truth—of the Divine, creation, and the Self—the \emph{sādhaka} attains ultimate liberation from \emph{saṁsāra}. 
\end{inspiration}
\newpage

\slcol{\Index{तेषामहं समुद्धर्ता} मृत्युंसंसारसागरात् ।\\
भवामि नचिरात्पार्थ मय्यावेशितचेतसाम् ॥ १२.७ ॥ }
\translate{teṣām ahaṁ samuddhartā mṛtyuṁ-saṁsāra-sāgarāt |\\
bhavāmi na cirāt pārtha mayy āveśita-cetasām || 12.7 ||}
\cquote{For those whose minds are absorbed in Me, O Pārtha (son of Pṛthā, Arjuna), I become the savior from the ocean of cycles of birth and death, without any delay.}
\slcol{\Index{मय्येव मन} आधत्स्व मयि बुद्धिं निवेशय ।\\
निवसिष्यसि मय्येव अत ऊर्ध्वं न संशयः ॥ १२.८ ॥}
\translate{mayy eva mana ādhatsva mayi buddhiṁ niveśaya |\\
nivasiṣyasi mayy eva ata ūrdhvaṁ na saṁśayaḥ || 12.8 ||}
\cquote{Fix your mind on Me alone, place your intellect in Me; then you shall dwell in Me alone, hereafter—of this there is no doubt.}
\slcol{\Index{अथ चित्तं समाधातुं} न शक्नोषि मयि स्थिरम् ।\\
अभ्यासयोगेन ततो मामिच्छाप्तुं धनञ्जय ॥ १२.९ ॥}
\translate{atha cittaṁ samādhātuṁ na śaknoṣi mayi sthiram |\\
abhyāsa-yogena tato mām icchāptuṁ dhanañjaya || 12.9 ||}
\cquote{But if you are unable to fix your mind steadily on Me, then O Dhanañjaya (the winner of wealth, Arjuna), through the system of repeated practice (\emph{abhyāsa-yoga}) seek to reach Me.}
\slcol{\Index{अभ्यासेऽप्यसमर्थोऽसि} मत्कर्मपरमो भव ।\\
मदर्थमपि कर्माणि कुर्वन् सिद्धिमवाप्स्यसि ॥ १२.१० ॥ }
\translate{abhyāse ’py asamartho ’si mat-karma-paramo bhava |\\
mad-artham api karmāṇi kurvan siddhim avāpsyasi ||12.10 ||}
\cquote{If you are unable even to follow such a system (\emph{abhyāsa-yoga}), then be intent on doing actions for My sake. By even performing actions for Me, you shall attain perfection.}

\newpage
\begin{mananam}{\mananamfont {Mananam Śloka - 6, 8}}
\normalfont\mananamtext Why do most people not prioritise God above all else? 
How does the Lord liberate me from \emph{saṁsāra}? Does this imply that divine intervention will instantly dissolve all of my problems?
What challenges prevent me from keeping my mind and intellect constantly focused on the Lord? What shifts in my understanding and behaviour are necessary in order to make God the central focus of my life?
\end{mananam}
\sreflect
\begin{inspiration}{\inspirationTitleFont{Inspiration}}
\normalfont\mananamtext For those entangled in the challenges and struggles of worldly life, the path of unwavering devotion to the Lord offers solace. When devotees embrace a cherished form of the Divine and engage in sincere devotional practice, they begin to perceive God permeating the essence of all worldly existence. Then, through the Lord’s grace, the heart is gradually purified, freeing it from outward temptations and inward delusions. A mind and intellect refined by such devotion become worthy instruments for realizing the ultimate truth of God and oneself.
\end{inspiration}
\newpage

\slcol{\Index{अथैतदप्यशक्तोऽसि} कर्तुं मद्योगमाश्रितः ।\\
सर्वकर्मफलत्यागं ततः कुरु यतात्मवान्  ॥ १२.११ ॥}
\translate{athaitad apy aśakto ’si kartuṁ mad-yogam āśritaḥ |\\
sarva-karma-phala-tyāgaṁ tataḥ kuru yatātmavān || 12.11 ||}
\cquote{But if you are unable even to do this, then, taking refuge in My \emph{yoga}, renounce the fruits of all actions with a self controlled mind.}
\slcol{\Index{श्रेयो हि ज्ञानमभ्यासा}ज्ज्ञानाद्ध्यानं विशिष्यते ।\\
ध्यानात्कर्मफलत्यागस्त्यागाच्छान्तिरनन्तरम्  ॥ १२.१२ ॥}
\translate{śreyo hi jñānam abhyāsāj jñānād dhyānaṁ viśiṣyate |\\
dhyānāt karma-phala-tyāgas tyāgāc chāntir anantaram ||12.12||}
\cquote{Better indeed is \emph{jñānam} (knowledge) than mere repeated practice (\emph{abhyāsa}); superior to  \emph{jñāna} is \emph{dhyānam} (meditation); superior to \emph{dhyānam} is renunciation of the fruits of action—for from such \emph{tyāga } (renunciation), peace follows immediately.}
\slcol{\Index{अद्वेष्टा सर्वभूतानां} मैत्रः करुण एव च ।\\
निर्ममो निरहङ्कारः समदुःखसुखः क्षमी ॥ १२.१३ ॥}
\translate{adveṣṭā sarva-bhūtānāṁ maitraḥ karuṇa eva ca |\\
nirmamo nirahaṅkāraḥ sama-duḥkha-sukhaḥ kṣamī || 12.13 ||}
\cquote{He who is free from hatred towards all beings, friendly, and compassionate, free from possessiveness and ego, equanimous in pleasure and pain, and forgiving....}
\slcol{\Index{सन्तुष्टः सततं योगी} यतात्मा दृढनिश्चयः ।\\
मय्यर्पितमनोबुद्धिर्यो मद्भक्तः स मे प्रियः ॥ १२.१४ ॥ }
\translate{santuṣṭaḥ satataṁ yogī \\\hspace*{0.5cm}yatātmā dṛḍha-niścayaḥ |\\
mayy arpita-mano-buddhir yo \\\hspace*{0.5cm}mad-bhaktaḥ sa me priyaḥ || 12.14 ||}
\cquote{Ever content, a self-controlled \emph{yogi}, firmly resolved, with mind and intellect dedicated to Me—such a devotee of Mine is dear to Me.}

\newpage
\begin{mananam}{\mananamfont {Mananam Śloka - 9, 12}}
\mananamtext\normalfont How can I cultivate the qualities described in these verses?
Among various spiritual practices, which one resonates most deeply with me?  
Do I feel more drawn to worship, meditation, or the study of the scriptures?  
In my actions, am I able to offer them all wholeheartedly to the Lord (\emph{arpaṇam})? Am I able to surrender expectation of specific results and accept both good and bad as the Lord’s blessings (\emph{prasādam})? 
How can I integrate different spiritual paths to support my growth, adapting to them as per my evolving needs and mental states?  
\end{mananam}
\sreflect
\begin{inspiration}{\inspirationTitleFont{Inspiration}}
\normalfont\mananamtext The scriptures present a wide range of spiritual practices and paths,
The scriptures not only present diverse spiritual paths to suit each \emph{sādhaka’s}
unique nature, but also offer flexible approaches that align with one's evolving mental states and life circumstances. As true devotee or seeker, we can never claim that our inner or outer conditions have limited our spiritual progress, as the sacred teachings continuously adapt to support our spiritual growth at every stage of the journey.
\end{inspiration}
\newpage

\begin{mananam}{\mananamfont {Mananam Śloka - 13, 14}}
\normalfont\mananamtext How does being free from hatred and embracing compassion for all make one beloved to the Lord? Why does the attachment to "I" and "mine" pull us away from devotion to the Lord?
In what way does maintaining equanimity in both joy and suffering help us to fix our mind on the Lord? How do self-restraint and unwavering conviction serve as pathways to attaining the Lord?
To what extent have I truly surrendered my mind and intellect to the Lord?
\end{mananam}
\sreflect
\begin{inspiration}{\inspirationTitleFont{Inspiration}}
\normalfont\mananamtext Just as a child longs for a mother’s love and a lover seeks the attention of her beloved, a devotee should strive to earn the Lord’s grace by cultivating qualities that are pleasing to Him. The virtues that endear one to the Lord require transcending afflictive traits such as egoism, hatred, and attachment. Through diligence and mindful effort, we as devotees gradually overcome these obstacles, refining the heart and purifying our intentions on the path to divinity.
\end{inspiration}
\newpage

\slcol{\Index{यस्मान्नोद्विजते लोको} लोकान्नोद्विजते च यः ।\\
हर्षामर्षभयोद्वेगैर्मुक्तो यः स च मे प्रियः ॥ १२.१५ ॥}
\translate{yasmān nodvijate loko lokān nodvijate ca yaḥ |\\
harṣāmarṣa-bhayodvegair mukto yaḥ sa ca me priyaḥ || 12.15 ||}
\cquote{He by whom the world is not disturbed, and who is not disturbed by the world, who is free from joy, intolerance, fear, and agitation—(such a devotee) is dear to Me.}
\slcol{\Index{अनपेक्षः शुचिर्दक्ष} उदासीनो गतव्यथः ।\\
सर्वारम्भपरित्यागी यो मद्भक्तः स मे प्रियः ॥ १२.१६ ॥ }
\translate{anapekṣaḥ śucir dakṣa udāsīno gata-vyathaḥ |\\
sarvārambha-parityāgī yo mad-bhaktaḥ sa me priyaḥ || 12.16 ||}
\cquote{He who is free from wants, pure, skillful, indifferent, free from agony, and who renounces all undertakings—such a devotee of Mine is dear to Me.}
\slcol{\Index{यो न हृष्यति न द्वेष्टि} न शोचति न काङ्क्षति ।\\
शुभाशुभपरित्यागी भक्तिमान्यः स मे प्रियः ॥ १२.१७ ॥}
\translate{yo na hṛṣyati na dveṣṭi na śocati na kāṅkṣati |\\
śubhāśubha-parityāgī bhaktimān yaḥ sa me priyaḥ || 12.17 ||}
\cquote{He who neither rejoices nor hates, neither grieves nor desires, who renounces both good and bad, and is full of devotion—(such a devotee) is dear to Me.}
\slcol{\Index{समः शत्रौ च मित्रे} च तथा मानापमानयोः ।\\
शीतोष्णसुखदुःखेषु समः सङ्गविवर्जितः ॥ १२.१८ ॥}
\translate{samaḥ śatrau ca mitre ca tathā mānāpamānayoḥ |\\
śītoṣṇa-sukha-duḥkheṣu samaḥ saṅga-vivarjitaḥ || 12.18 ||}
\cquote{He who remains same towards friend and foe, and also in honour and dishonour; who is same in cold and heat, pleasure and pain, and free from attachment…}

\newpage
\begin{mananam}{\mananamfont {Mananam Śloka - 15}}
\normalfont\mananamtext Either in my thoughts or through deeds, do I cause agitation or negativity to others? How are people affected by my presence or through my speech and actions?
Am I easily affected by others’ presence and energy? Am I too sensitive and touchy? Am I envious and anxious of certain people? Am I easily hurt by others’ criticism and attitude? What shift in understanding or attitude can give me the strength to stand unshaken internally despite all challenges externally?
\end{mananam}
\sreflect
\begin{inspiration}{\inspirationTitleFont{Inspiration}}
\normalfont\mananamtext A true devotee of God should refrain from causing unnecessary harm or disruption to others, except in rare cases where fulfilling one’s righteous duties makes it unavoidable. Just as we seek peace and harmony in our own lives, we should also strive to cultivate and share these qualities with those around us.
However, being considerate of others should not weaken our inner strength or make us overly affected by external influences. Instead, we must develop a resilient mind, firmly rooted in faith in the Lord’s justice and compassion, allowing us to navigate life with both sensitivity and steadfastness.
\end{inspiration}
\newpage

\begin{mananam}{\mananamfont {Mananam Śloka - 16}}
\normalfont\mananamtext Why is contentment essential in receiving the Lord’s grace? How do personal desires lead to agitation and restlessness? 
What distinguishes personal desires (\emph{sakāma}) from impersonal desires (\emph{niṣkāma})? How can I integrate the well-being of others into my own personal aspirations? 
As a devotee of God, what should my true motivation be while undertaking a new project or task? In what ways can ambitious pursuits—requiring time, effort, and resources—distract me from my focus on God?
\end{mananam}
\sreflect
\begin{inspiration}{\inspirationTitleFont{Inspiration}}
\normalfont\mananamtext When a desire takes root in the heart, it stirs restlessness, inevitably leading to action—whether in thought or deed. While action itself is not wrong, only those deeds that are dedicated to the well-being of others and offered to the Lord are truly purifying. Such selfless actions, \emph{niṣkāma} \emph{karma}, bring liberation, whereas actions driven by personal gratification, \emph{sakāma} \emph{karma}, create attachment and bondage.
For one bound by personal desires, peace remains elusive, as one longing leads to another in an unending cycle. But the devotee who acts solely for the Lord, unattached to success or failure, moves through life with inner contentment, fulfilling all duties without disturbance.
\end{inspiration}
\newpage

\begin{mananam}{\mananamfont {Mananam Śloka - 17, 18}}
\normalfont\mananamtext How do the extremes of dualistic worldly experiences—the pairs of opposites—affect me?
Why is it essential to cultivate detachment from both success and failure, praise and criticism, harmonious relationships and conflicting ones?
The \emph{Gītā} advocates mental equanimity. Does this serve as a means of emotional escapism (a mere coping mechanism), or does it hold a deeper significance in terms of spiritual progress?
\end{mananam}
\sreflect
\begin{inspiration}{\inspirationTitleFont{Inspiration}}
\normalfont\mananamtext Equanimity is a definitive mark of spiritual progress. Generally, our attention tends to flow outwards—to objects, people, and the external world—while the mind habitually labels experiences as favourable or unfavourable.
A true spiritual practitioner must cultivate the ability to remain present with every experience in a state of pure observation. At an advanced stage, this includes the skill of witnessing even one’s own internal feelings and judgements, without involvement. Mastering this supreme art of detached awareness is to realise the divine presence within us.
\end{inspiration}
\newpage

\slcol{\Index{तुल्यनिन्दास्तुतिर्मौनी} सन्तुष्टो येन केनचित् ।\\
अनिकेतः स्थिरमतिर्भक्तिमान्मे प्रियः नरः ॥ १२.१९ ॥}
\translate{tulya-nindā-stutir maunī santuṣṭo yena kenacit |\\
aniketaḥ sthira-matir bhaktimān me priyaḥ naraḥ || 12.19 ||}
\cquote{He who is the same in blame and praise, silent, content with anything, homeless, steady in mind, and full of devotion—that person is dear to Me.}
\slcol{\Index{ये तु धर्म्यामृतमिदं} यथोक्तं पर्युपासते ।\\
श्रद्दधाना मत्परमा भक्तास्तेऽतीव मे प्रियाः ॥ १२.२० ॥}
\translate{ye tu dharmyāmṛtam idaṁ yathoktaṁ paryupāsate |\\
śraddadhānā mat-paramā bhaktās te ’tīva me priyāḥ || 12.20 ||}
\cquote{Indeed those devotees who follow this immortal \emph{dharma} as described above, endowed with faith, regarding Me as their supreme goal—they are extremely dear to Me.}

\newpage
\begin{mananam}{\mananamfont {Mananam Śloka - 19, 20}}
\normalfont\mananamtext How does devotion to the Lord bring inner contentment and allow me to rely on Him as my ultimate refuge for all needs? 
Why does the Lord bestow special grace upon those who make Him the supreme goal of their lives? Can I still receive the Lord’s grace if I pursue other goals alongside Him?
Why is it important to follow the \emph{dharma} prescribed by the Lord? Is devotion alone sufficient, or must I also align my actions with His teachings?
\end{mananam}
\sreflect
\begin{inspiration}{\inspirationTitleFont{Inspiration}}
\normalfont\mananamtext When we deeply love and respect someone, we naturally seek to follow that person’s guidance and fulfill his or her expectations. True devotion to the Lord involves trusting His divine laws and wholeheartedly aligning with the eternal \emph{dharma} He has established for our well-being—passed down through the scriptures and the sacred lineage of \emph{ṛṣis}.
By dedicating ourselves to the Lord and walking the path of righteousness, no injustice against us will remain unresolved. The omniscient and all-knowing Lord ensures that we are protected and that our essential needs are always met through His grace. 
\end{inspiration}
\newpage

\begin{center}
\devfont ॐ तत्सदिति श्रीमद्भगवद्नीतासूपनिषत्सु  ब्रह्मविद्यायां  \\योगशास्त्रे
श्रीकृष्णार्जुनसंवादे  भक्ति योगो नाम द्वादशोऽध्यायः॥\\
\chapEndSlokaTranslate{ōṃ tatsaditi  śrīmadbhagavadnītās ūpaniṣatsu\\  brahmavidyāyāṁ  yōgaśāstrē\\
śrīkṛṣṇārjunasaṁvāde  bhakti yōgō nāma dvādaśō ’dhyāyaḥ ||}
\end{center}
